% ==================================================
% LEMBAR PENGESAHAN
% ==================================================

\clearpage

\phantomsection
\addcontentsline{toc}{chapter}{LEMBAR PENGESAHAN}


\begin{center}
    {\fontsize{14pt}{16pt}\selectfont\bfseries
    \MakeUppercase{LEMBAR PENGESAHAN}\\
    \MakeUppercase{PROPOSAL DISERTASI}}
\end{center}

\vspace{0.7cm}

\noindent
{\fontsize{12pt}{17pt}\selectfont}
\begin{tabular}{@{}p{2cm}@{}p{0.4cm}@{}p{11.6cm}@{}}
    Judul & : & \varJudulID \\
    Oleh & : & \varMhsNama \\
    NRP  & : & \varMhsNrp \\
\end{tabular}

\vspace{0.5cm}

\begin{center}
    {\fontsize{14pt}{16pt}\selectfont\bfseries
    {Telah diseminarkan pada :}}
\end{center}

\vspace{0.1cm}

\noindent
{\fontsize{12pt}{17pt}\selectfont}
\begin{tabular}{@{}p{2cm}@{}p{0.4cm}@{}p{11.6cm}@{}}
    Hari & : & \varUjianKualifikasiHari \\
    Tanggal & : & \varUjianKualifikasiTanggal \\
    Tempat  & : & \varUjianKualifikasiRuangTempat \\
\end{tabular}

\vspace{0.5cm}

\begin{center}
    {\fontsize{12pt}{16pt}\selectfont
    {Mengetahui / menyetujui:}}
\end{center}

\begin{table}[H]
    \renewcommand{\arraystretch}{1.15} % Ini mengatur tinggi baris tabel. Nilai default kira-kira 1.0.
    \setlength{\tabcolsep}{3pt} % mengatur jarak horizontal antara isi sel dengan garis/batas kolom.

    \begin{tabular}{
        >{\raggedright\arraybackslash}m{0.2cm}
        >{\raggedright\arraybackslash}m{7cm}
        >{\raggedright\arraybackslash}m{0.2cm}
        >{\raggedright\arraybackslash}m{7cm}
    }
        % \hline
        
        % \multicolumn{1}{|>{\centering\arraybackslash}m{4.3cm}|}
        {\textbf{}} &
        \textbf{Dosen Penguji :} &
        \textbf{} &
        \textbf{Dosen Pembimbing :} \\[3cm]
        % \hline

        1. &
        \varPengujiSatu &
        1. &
        \varPromotorSatu \\
        % \hline

        &
        NIP.~\varPengujiSatuNIP&
        &
        NIP.~\varPromotorSatuNIP \\[3cm]
        % \hline

        2. &
        \varPengujiDua &
        2. &
        \varPromotorDua \\
        % \hline

        &
        NIP.~\varPengujiDuaNIP&
        &
        NIP.~\varPromotorDuaNIP \\
        % \hline
    \end{tabular}

\end{table}

\clearpage

